\section{总结与延伸}

本期深入剖析了 LightRAG 的四个核心技术细节：

\begin{enumerate}
    \item \textbf{KG-based RAG 的必要性}：解决了 Naive RAG 在多跳推理和全局查询上的根本局限
    \item \textbf{增量索引机制}：通过 MapReduce 式的描述合并和语义压缩，支持动态新增文档
    \item \textbf{KG 可视化}：基于 pyvis 提供直观的图谱浏览工具
    \item \textbf{多模式检索}：Naive/Hybrid/Mixed 三种模式覆盖不同的查询需求
\end{enumerate}

\subsection{拓展阅读}

\begin{itemize}
    \item LightRAG 源码中的 \texttt{prompts.py}：包含所有关键 Prompt（Entity Extraction、Summarization 等）
    \item LangFuse 本地部署：用于监控 LightRAG 所有 API 调用的输入输出
    \item Microsoft GraphRAG：与 LightRAG 类似的 KG-based RAG 方案，可做对比研究
\end{itemize}

%% --- Fin del contenido --- %%

\end{document}
